\section{第一层塌陷：功能性危机}

存在、凭据、准入、承受，前面两节讨论的都是结构；结构的变化要落到人身上，才会显示出它真正的代价。这一步的落点不在收入，也不在岗位，而在一个更早的位置：一个人原来用什么证明自己在这个系统中占有一个位置。这一层的变化发生得最快，也最容易被误读成个人问题。
\begin{figure}[ht]
\centering
% >>>>>>>>>> figs/fig-collapse-spiral.tex >>>>>>>>>>
% 图 5-1  三层塌陷：变化速度依次放慢，察觉难度依次上升
% 三层不是三个并列的问题，而是一条因果链：功能层的真空把人推向别处，
% 别处提供的是一组保留刺激、取消做功的系统，在这种环境里边界不再被使用。
% 左侧色块给出每层的变化速度与察觉难度，右侧盒子给出机制与落点。
% 三层的颜色由 terra 经 rust 到 forest，只表示深度递增，不表示价值高低。
\begin{tikzpicture}[x=1cm, y=1cm, font=\small,
  bx/.style={rounded corners=2pt, line width=0.7pt},
  st/.style={rounded corners=1.5pt, line width=0.5pt},
  tt/.style={font=\small\heiti, anchor=west},
  tx/.style={font=\scriptsize, text=ink, anchor=north west, align=left},
  sl/.style={font=\scriptsize, text=slate, anchor=center, align=center},
  cy/.style={font=\scriptsize, text=ink, anchor=west}
]
\useasboundingbox (-0.20,0.40) rectangle (13.15,-7.05);

\def\sA{0.15} \def\sB{2.55} \def\bA{2.85} \def\bB{12.95} \def\arX{3.10}

% ================= 第一层：功能性危机 =================
\fill[terra!13] (\bA,0.00) rectangle (\bB,-1.85);
\draw[bx, terra!60] (\bA,0.00) rectangle (\bB,-1.85);
\fill[terra!26] (\sA,0.00) rectangle (\sB,-1.85);
\draw[st, terra!60] (\sA,0.00) rectangle (\sB,-1.85);
\node[tt, text=terra!95] at ({\bA+0.16},-0.34) {第一层塌陷｜功能性危机};
\node[tx] at ({\bA+0.16},-0.62)
      {接口失效，不是能力消失：会写、会算、会总结、会表达、会计划，\\[1pt]
       这些用来证明位置的能力被系统接手。落点是人原来凭什么证明自己占有一个位置。};
\node[sl] at ({\sA+1.20},-0.72) {变化最快};
\node[sl] at ({\sA+1.20},-1.10) {最易被误读};
\node[sl] at ({\sA+1.20},-1.44) {为个人问题};

% ---------- 层间传递 ----------
\draw[ink!60, line width=0.7pt, -{Stealth[length=1.8mm]}] (\arX,-1.87) -- (\arX,-2.43);
\node[cy] at ({\arX+0.15},-2.15) {功能层的真空把人推向别处};

% ================= 第二层：真实性危机 =================
\fill[rust!13] (\bA,-2.45) rectangle (\bB,-4.30);
\draw[bx, rust!60] (\bA,-2.45) rectangle (\bB,-4.30);
\fill[rust!26] (\sA,-2.45) rectangle (\sB,-4.30);
\draw[st, rust!60] (\sA,-2.45) rectangle (\sB,-4.30);
\node[tt, text=rust!95] at ({\bA+0.16},-2.79) {第二层塌陷｜真实性危机};
\node[tx] at ({\bA+0.16},-3.07)
      {系统提供结果，过程被吸收：输入几个词就得到一幅画，发出请求就得到回应。\\[1pt]
       变化是刺激与后果解耦——保留刺激，取消做功，完成感不再需要付出代价。};
\node[sl] at ({\sA+1.20},-3.17) {变化居中};
\node[sl] at ({\sA+1.20},-3.55) {表现为};
\node[sl] at ({\sA+1.20},-3.89) {「更便利」};

% ---------- 层间传递 ----------
\draw[ink!60, line width=0.7pt, -{Stealth[length=1.8mm]}] (\arX,-4.32) -- (\arX,-4.88);
\node[cy] at ({\arX+0.15},-4.60) {刺激与后果解耦之后，拒绝失去对象};

% ================= 第三层：主体性危机 =================
\fill[forest!13] (\bA,-4.90) rectangle (\bB,-6.75);
\draw[bx, forest!60] (\bA,-4.90) rectangle (\bB,-6.75);
\fill[forest!26] (\sA,-4.90) rectangle (\sB,-6.75);
\draw[st, forest!60] (\sA,-4.90) rectangle (\sB,-6.75);
\node[tt, text=forest!95] at ({\bA+0.16},-5.24) {第三层塌陷｜主体性危机};
\node[tx] at ({\bA+0.16},-5.52)
      {边界不是被夺走的，而是在没有阻力的环境里不再被使用。\\[1pt]
       系统不再需要你说「不」：拒绝失去对象，判断因此没有校准的机会。};
\node[sl] at ({\sA+1.20},-5.62) {变化最慢};
\node[sl] at ({\sA+1.20},-6.00) {最难在个体};
\node[sl] at ({\sA+1.20},-6.34) {时间尺度察觉};
\end{tikzpicture}

% <<<<<<<<<< figs/fig-collapse-spiral.tex <<<<<<<<<<

\caption{三层塌陷的传递关系。左侧色块给出三层的性质，右侧方框是各自的机制，层间箭头给出由此及彼的传导方向。三层的察觉难度依次上升，变化速度则互不同步；本章以下三节按这条链的顺序展开。}
\label{fig:collapse-spiral}
\end{figure}


\subsection{功能性失重：接口失效不是能力消失}

\subhead{人用它来证明位置的那些能力}

前面几节的结论落在承受能力上，但承受能力成为价值源点，并不意味着旧的价值证明方式会平稳过渡到新的方式上。过渡期间先发生的事情是失效。在过去的几十年里，一个人在社会中的位置，很大程度上由一组能力支撑。会写，意味着能把想法整理成可以被他人使用的形式；会算，意味着能在模糊信息中给出可依赖的数字判断；会总结，意味着能从大量材料中提取结构；会表达，意味着能在会议、谈判、课堂里稳住局面；会计划，意味着能把目标拆成步骤、资源和时间。

这组能力有一个共同点：它们都是可交付的。交付物可以被检验，可以进入流程，可以产生经济价值。因此具备这些能力的人不需要额外说明自己的价值，交付物本身就是证明。这也意味着，这组能力与人的社会位置之间存在一条直接的绑定：一个人会什么，决定了他能进入什么位置，能在多少事情上被需要，能被谁记住。这条绑定如此稳定，以至于它常常被当成人的价值的定义，而不只是价值的一种表现形式。生成式系统改变的正是这条绑定。

\subhead{接口失效不是能力消失}

需要把变化的位置说准确，因为一种流行的说法把它说错了地方。问题不在于这些能力变得无用，而在于它们不再具有分辨力。当一份结构清晰的报告可以由系统在几十秒内生成，写得清楚就不再能说明什么；当数据分析的基本流程可以被自动完成，算得准确就不再能区分谁受过训练。这不是因为人写得不好或者算得不对，而是因为产出这个尺度本身的信号强度大幅下降了，而信号强度的下降与能力水平的高低是两件事。\key{AI 不是立刻取消人的价值，而是取消许多人证明价值的旧接口。}

把分辨力写成一个判别力指标，

\begin{equation}
d'=\frac{\Delta\mu}{\sigma}
\end{equation}

式中 $\Delta\mu$ 是在同一评价尺度上两类人的产出均值之差，例如受过训练者与未受训练者的产出质量均值之差；$\sigma$ 是该尺度上个体产出的标准差；二者必须在同一尺度和同一人群中估计，相除后 $d'$ 无量纲，表示这个尺度把两类人区分开的能力。这个式子的含义是：$\Delta\mu$ 的收缩与 $\sigma$ 的不收缩可以同时发生。当合格产出变成默认可得，均值差被系统压平，而由生成引入的波动并不随之下降，$d'$ 因此趋近于零，尺度失去分辨力，而两类人的绝对水平都没有改变。接口失效不是能力消失，写出来就是失效发生在 $d'$ 上，不发生在 $\mu$ 上。适用条件是尺度可比且两组方差大致相当，方差相差很大时判别力需要换一种定义。它刻画的是尺度作为凭据的鉴别力，不是任何人的能力值，本身不给出实测数值。

这里需要把接口失效和失业区分开，因为混同两者会同时高估和低估问题。失业是一个岗位问题，接口失效是一个证明问题。许多人的岗位还在，因为他们的工作处在高黏滞环节：需要现场、需要身体、需要长期关系、需要承担后果；但他们的自我评价方式已经被改组了。一个人依然每天在岗位上完成工作，可他不再确定自己的位置是因为“我会这个”。他会的东西仍在，但不再稀缺；他仍然被需要，但被需要的原因变成了别的什么。这种落差在具体生活中的表现相当一致：一个在自己的领域里做了二十年的人，专业判断已经内化为直觉，现在他打开一个系统，输入几句话，得到的输出覆盖了他百分之八十的日常工作；他仍然要判断，仍然要挑选，仍然要对结果负责，但工作的重心从生产前移到了审核，而审核在自我认同上不提供与生产同等的支撑。

\subhead{劳动、制作与行动}

这个落差之所以深，是因为它触及的不只是收入，而是活动与承认之间的关系，而这一层关系已经有人做过细致区分。Hannah Arendt 在分析现代社会时把人的活动分成三种：劳动，对应生物性的新陈代谢，产物迅速被消耗；制作，对应制造耐久之物，产物比制造者活得长；行动，对应在他人面前说话和做事，不产生物件，却创造关系与历史。\cite{arendt1958human} 她的判断是，现代社会把这些活动都还原成了劳动，一切以生命过程的维持为尺度来衡量，制作和行动失去了独立地位。

这个区分可以用来定位当前的情况。旧的能力接口之所以稳固，是因为它同时连接着制作和行动：一个人会做某件事，既生产了耐久的东西，也在共同体里获得了一个位置和一段历史，产出既是他与物的关系，也是他与人的关系。生成系统接管的主要是制作环节，而且是其中可标准化、可复制的那部分。当制作被接管，与它捆绑在一起的那部分社会承认也跟着松脱。这不是因为社会承认被剥夺了，而是因为它的载体变了，而新的载体还没有建立起来。

分歧需要在这里说清楚，否则Arendt只是被借来做一个诊断，而不是被当作对手。她的方案落在行动上：行动的场域是公共领域，价值不在产物，而在复数的人彼此可见、彼此应答，一个人由此显现为一个“谁”而不是一个“什么”\cite{arendt1958human}。本书给不出这个方案，它拿出的替代品是承受：一个位置值得被占据，是因为后果回到那里。两者不是同一个东西。行动的价值是内在的，它要求复数性，要求有他人看见并回应；承受的价值是功能的，它只要求后果确实回流。按本书的判据，一个独自承受全部后果的存在物完全合格，而Arendt不会承认这种存在物有资格进入她的公共领域。

这个替换有代价，代价要付在明处。以承受为准，本书能够指出后果回到哪里，却不能说明回到那里的负担是否应当由那一方承担。判据本身对分配是沉默的，在最坏的情况下它会把任何既成的负荷都读成价值的证据。Arendt的政治语言保住了这一层：她至少可以问这是不是不义。本书把失去功能位置这件事翻译成结构语言之后，得到的是一份更精确的地图，失去的是在地图上标注正当性的那支笔。

Arendt在另一处用过一个概念，指称一个群体在社会结构中失去功能位置的状态，她讨论的是特定历史条件下的政治后果。本章借用它的结构含义：当系统状态转移不再需要某个节点参与时，这个节点并没有被消灭，它只是从流程中滑了出来。\key{滑出的位置没有自动的替代品，它是空的。}这个说法容易被读成一种悲观的断言，但它实际上是一个关于自动性的说明：市场和制度都不会自动为滑出的节点安排新的功能位置，因为没有人有动机去做这件事：对每一个具体的组织而言，那个节点的缺席并不构成问题。

\subhead{为什么这个诊断不能被简化成情绪问题}

正因为位置不会自动补上，关于这种失重的流行解释才格外值得检查，而它们恰好都绕开了这一点。一种是安慰性的：人的能力并没有消失，只是工具变了，学会用工具就好。这个说法在操作层面成立，在结构层面回避了问题。学会使用工具之后，人与产出之间的绑定关系仍然没有恢复，因为工具放大的正是可以自动化的那部分能力，而放大的结果是这部分能力更加不能作为区分依据。另一种是责备性的：用不了新工具的人是自己不努力。这个说法混淆了两件事，个体层面的适应能力和结构层面的接口重置，而再努力的人也无法通过提升自己来解决一个已经不依赖该能力的系统问题。

两种解释共同的错误，是把问题定位在个体身上，而实际情况是被置换的是接口，不是人。接口这个词值得强调，因为它决定了应对的方向：接口是一组约定的匹配关系，这边提供什么，那边需要什么，两者如何对接；接口失效意味着匹配关系不再成立，而不意味着任何一方有缺陷。当一批接口同时失效，受影响的人数以千万计，那么它就是结构问题，而结构问题的解决位置在制度一侧，不在个人一侧。把这个判断说清楚不是为了免除个人责任，而是为了指认责任的正确位置。

\subhead{诊断的边界}

这个诊断需要划清界限，否则会滑向两个它并不支持的结论，而这两个方向都会让整章论证失去分量。它不主张任何具体职业会消失。接口失效与岗位取消是两件事，前者的发生范围远大于后者，而且发生得更早；把接口失效等同于失业预测，会低估问题的普遍性，因为它影响的是大量仍然在职的人。它也不主张专业训练没有意义。在高后果领域，训练的作用不是生产可交付物，而是形成判断力：一个受过训练的医生与一个没受过训练的人使用同样的系统，得到结果的质量差别很大，区别不在于谁知道得更多，而在于谁能在结果出错时识别出来。这种能力不会因为生成成本下降而贬值，反而会升值。

真正被削弱的是中间那一层：不需要深度判断、也不承载高风险后果，只需要完成标准化制作的那一层。这一层在过去是最大量的就业，也是最多人建立职业认同的地方，因此它的收缩影响的人数最多，而它得到的关注最少，因为它既不属于需要保护的弱势群体，也不属于需要争夺的高端人才。失重发生在这里，代价也落在这里。而它引发的第一反应通常不是抵抗，而是退向一个更容易获得完成感的地方，这与前面讨论过的低黏滞路径正好接上。退向完成感只是结果，原因在更早的地方：在那些被压缩掉的环节里，人失去的不只是交付物，还有从确立目标到克服阻力那条回路本身。

\subsection{权力的断裂与根本性的多余}

接口失效说明的是位置的变化，它还没有解释这种变化为什么会落到一个人的自我评价上。一个人在岗位上仍然被需要，却不再确定自己为什么被需要，这中间隔着一个更早的环节：从确立目标，到克服阻力，再到得到结果的那条回路。这条回路不通过产出，而通过做功本身，把人绑在世界上。

\subhead{权力过程与回路短路}

Theodore Kaczynski 在《工业社会及其未来》里给出过一个解释。他把人的满足感追溯到一条他称为“权力过程”的回路：确立一个目标，为它付出努力，克服途中的阻力，最后达成它；满足感不来自结果的获得，而来自阻力被实际克服的那一段\cite{kaczynski1995industrial}。这个描述不必连同他的作者身份一起接受，也不必连同他对技术文明的结论一起接受。作为一个机制描述，它与本书前面的分析是同一个结构：目标是一个待实现的状态，阻力是达成它的代价，做功是支付代价的过程。工业技术系统改变的是这条回路的中段。系统把结果直接交付给使用者，把克服阻力的部分留给自己。人得到成品，失去做功。

在 RTG 的语言里，这一步可以写得更准确。当外部系统替人完成一次状态转移时，这条转移对人而言的 R 趋近于零；T 被压缩到几乎不可感知；G 也不再需要人的参与。三个维度同时归零，意味着这件事对人不再构成一次做功，也就不再产生与做功相配的反馈。写成极限，
\begin{equation}
\eta_{p}=\eta_{0}\exp(\alpha R_{p}+\beta T_{p}+\gamma G_{p}),
\qquad
\lim_{(R_{p},\,T_{p},\,G_{p})\to(0,0,0)}\eta_{p}=\eta_{0}
\label{eq:trip-zero}
\end{equation}
式中 $\eta_{p}$ 是这一次状态转移对当事人而言的黏滞，下标 $p$ 表示各量都按当事人这一侧结算，而不是按系统的物理账本结算；$R_{p}$、$T_{p}$、$G_{p}$ 是同一件事在当事人坐标里的资源势垒、时间刚性与拓扑耦合；$\eta_{0}$ 是三个坐标都取零点时的参照黏滞。这个式子的要点不在极限值，而在它把三件事合并成一个：$\eta_{p}$ 退回 $\eta_{0}$，意味着这次转移对当事人既不耗费他掌握的什么，也不占用他的时间，也不经过他与任何他人的连接。三项同时消失，而不是各自少一点。与做功相配的反馈量正比于 $\eta_{p}-\eta_{0}$，所以当它退到零，反馈也一并退到零。这就是“人获得的是结果，不是完成”的算式形式。这里要防止把结论推得过宽：少量、偶发的结果交付不会改变什么，问题出在回路的大部分环节都被外部接管之后，剩下的环节不足以支撑一个人的位置感。

回路被接管之后，通常会出现两种补偿。一种是转向他处：在系统不关心的领域里自己设定目标、自己设定阻力，竞技、收藏、圈层内的地位竞争都属于这一类。这些活动在形式上是完整的，有目标，有阻力，有结果；它们与真实目标的差别不在强度，而在条件：目标可以自己修改，阻力可以自己降低，因此反馈携带的信息量低于真实目标，而信息量的差别决定了它能不能用来校准自己。另一种补偿是降低期待：不再设定需要长期做功的目标，把满足感的来源移到即时的、低代价的刺激上。

第二种补偿在心理学上有一个更早的名字。Martin Seligman 和 Steven Maier在上世纪六十年代的实验中让动物反复经历无法通过自身行为改变的伤害，随后发现，即使后来出现了可以逃脱的条件，动物也不再尝试\cite{seligman1967failure}。这一解释在五十年后被原作者修正过：被动的默认反应并不是学来的，学到的其实是“自己的行为能否影响结果”这一判断\cite{maier2016learned}。修正之后的结果对本章更有用：问题不在于人变得被动，而在于行为与结果之间的连接是否存在。连接在，人就会尝试；连接断掉，尝试就不再发生。这与回路短路是同一件事的两种说法。

Kaczynski对这套系统的判断比回路描述更重。他认为工业技术系统一旦形成，就必须持续扩张才能维持自身，它需要不断吸收新的能量、新的资源和新的解决方案；停止扩张不是稳态，而是解体，而个体对系统的依赖一旦形成，退回旧方式就不再是选项。这个判断有时被概括成对这个系统的热力学诅咒：它不能停，也不能退。这里需要划一条界线。可以保留的是回路结构，不能保留的是关于终局的预言。扩张是否必然，依赖是否不可逆，取决于制度与选择，而这两样在他的论证里被当成了常量。

分歧还不止这一处。把技术当作一个单一的行动者，会让黏滞的分布看起来是技术的固有属性；而在本书的框架里，摩擦削减到哪里、位置如何重组，取决于产权、劳动制度、公共投入和谈判能力。铁路、电气化和流水线都曾削掉大量旧摩擦，它们在英美、德国和日本造出的位置结构却各不相同：同样的机器，在不同的制度里替换掉的是不同的人，留下的也是不同的位置。这一点直接决定出口的形状。他的出口是摧毁工业体系，而摧毁恰好取消了本书要做的那项区分：哪些摩擦是无效的、应当消除，哪些摩擦承担着承付功能、必须保留。把两类摩擦一起消掉，得到的不是自由，而是更彻底的多余。

\subhead{根本性的多余}

回路短路解释的是“不需要我动手”，它还没有解释“不需要我”。后一个问题在另一个方向上早已被描述过。Arendt 在《极权主义的起源》里把“多余”当作一个结构范畴使用，而不是当作贫困的同义词：多余描述的是一种位置，没有人需要他的行动，他的存在不改变任何结果的分布；她进一步指出，由多余的人组成的群众最容易接受一种把一切事件解释得完全一致、并且不允许例外的系统\cite{arendt1951origins}。

把这个判断换成物理语言，可以得到更清楚的形式。在旧结构里，人是一种负熵的贡献者：通过劳动和判断，把外部低熵资源转化为可以维持的结构，系统因此需要他参与能量与信息的交换。当系统不再需要他参与这个交换，他的功能位置就断开了。断开不等于被剥夺。剥削的前提仍然是人有用，只是所得与付出不成比例；多余比剥削更冷，因为它不需要从人身上取得任何东西。\key{剥削仍然需要人是有用的，多余意味着连利用都不再需要；功能性危机最冷的部分，不是位置被夺走，而是位置本身不再存在。}这种失重很难被政治语言接住：抗议需要有具体的要求，而这里没有一个可以要求返还的东西。

多余还有一个麻烦之处，它在统计上几乎不可见。失业有口径，失业率、领取救济的人数、岗位空缺都可以计数，因此失业能够进入议程；多余没有口径，因为它常常发生在仍然在职的人身上。一个人可以在岗而多余：他的产出里由系统生成的部分不断增加，剩下的工作是监督系统、修补系统的输出、承接系统不愿承担的例外，也就是吸收残差。这类位置在统计上表现为就业，在结构上表现为缓冲层。它的稳定性低于它看上去的样子，因为一旦系统把残差也接管了，这个位置就不再有任何内容需要人参与。

这个诊断同样需要防止两种简化。一种是把多余写成命运的判决。功能位置的消失是可逆的，历史上有大量节点从流程中滑出之后，又在别处被重新需要，只是新位置通常与旧位置不在同一个地方。这正是后面几节要处理的问题。另一种是把责任推到某个具体主体身上。位置的消失不是阴谋的结果，它是一连串局部优化的叠加：每一个减少环节的决策在它自己的尺度上都合理，叠加起来才产生出一个没有人需要的节点。指认结构不等于免除任何人，而是为了把责任的落点放对位置。

权力的断裂和位置的消失合起来，解释了功能性危机为什么同时表现为两件事：一端是“不需要我动手”，一端是“不需要我”。两端之间的地带最宽，也最难被看见。大多数人既没有被完全替代，也没有被完全需要，他们所在的那一层正在被压缩，而压缩的过程不产生任何可以被抗议的事件。

\subsection{可被压缩的平庸与真空}

权力断裂的两端是少数人的处境：被完全接管的一极和被完全需要的一极。大多数人落在中间，而中间这一层的变化比两端更难描述，因为它既不以失业的形式出现，也不以失重的形式出现。它表现为一种压缩：需求被满足得越来越便宜，而满足需求的位置越来越不需要特定的人。

\subhead{被压缩的需求与位置的真空}

把需求层级换成物理语言，可以看清压缩发生在哪一层。常见的读法把安全、归属、尊重、自我实现理解为一串心理状态，但每一层都对应一种与外部结构的连接方式：底层的对应是能量输入，也就是收入与物资；中间层的对应是拓扑位置，也就是你在这个系统中占有的那条边；上层的对应是可干预范围，也就是你的判断能改变什么结果。三个连接的物理性质不同，受压的方式也不同。

机器智能的压缩首先落在中间层。可以被标准化的制作环节，包括整理、转写、摘要、初级编码、常规检索和模板化设计，恰好构成中层位置的绝大部分内容。压缩发生在供给侧：当生成这些产物的成本下降几个量级，产物本身不再稀缺，附带在产物上的位置也随之失去支撑。技能的通胀指的就是这件事。会写、会算、会总结这些能力在名义上仍然存在，但它们的存量相对于需求已经过剩，因此不再能作为区分人的依据。它与货币通胀有一点相同：名义数量上升，实际购买力下降，而持有者是最晚察觉的人。

这个机制有一个更早的对应物，它发生在学历上。教育扩张并不总是由技能需求驱动的：当高学历的供给增加，而对应职位不按同比例增加时，学历的作用就从技能证明转变为筛选门槛，个人被迫为相同的职位投入更长的教育时间。Randall Collins 在 1979 年把这一机制称为凭据社会，它的推论是，教育回报的一部分只是位置竞争的结果，而不是生产率的提高\cite{collins1979credential}。生成式系统把同一机制推进到技能层：当一种技能的供给可以按需扩张，它作为凭据的功能就会贬值，而贬值最先打击的是那些以该技能为唯一凭据的人。凭据通胀与货币通胀还有一个共同点，它不降低任何人的实际能力，只降低能力的兑换率。

底层的连接受冲击的方式不同。需要现场、身体、即时反应和持续在场的环节，介质不同，成本结构不同，短期内不容易被替代，也不容易被抬价。上层的位置在部分领域反而升值，因为高后果领域的判断必须有人承担。被压缩得最严重的是中间那一层，而这一层在过去是最大量的就业所在，也是最多人建立职业认同的地方。

任务替代与岗位变化之间，可以用ATM的历史区分。Bessen讨论过自动取款机减少部分柜面任务，同时分支机构和客户服务的变化创造其他需求。\cite{bessen2015learning} 这个案例显示的是任务、岗位与组织共同变化的机制，不能据此保证每一轮自动化都产生同样结果。

\begin{table}[ht]
\centering
\caption{ATM案例中的任务与组织变化}
\label{tab:ch5-atm}
{\small\begin{tabularx}{\textwidth}{@{}lXX@{}}
\toprule
层面 & 自动化的直接作用 & 还须观察的变化\\
\midrule
任务 & 取款与部分柜面操作被机器承担 & 每个网点的人工需求\\
组织 & 分支和服务方式可能调整 & 网点数量与服务范围\\
岗位 & 客户交互及其他职责重新组合 & 总就业、培训与薪酬\\
\bottomrule
\end{tabularx}}
\end{table}

表\ref{tab:ch5-atm}区分两种变化：机器接管一部分可标准化任务，组织和服务需求改变剩余岗位。2026年8月修订的斯坦福研究使用截至6月的工资记录，没有发现广泛的全经济岗位替代，同时记录了年轻劳动者在高暴露职业中的就业分化，并强调这是描述性指标。\cite{brynjolfsson2026canaries} 总量与入口机会，需要分别观察。

被压缩的只是位置，另一侧仍然存在大量没有被满足的需求，这个不对称值得单独说明。照护、教育、维修、社区服务这类环节的需求随人口结构与生活水平持续上升，但它们能提供的位置数量受预算、许可制度和财政能力约束，不随需求自动增加。这些领域还有一个共同点：产出难以标准化，单位成本不随时间下降，因此它们的相对价格持续上升，可雇佣的数量进一步受限。需求在而预算不在，结果是空位与缺人同时存在：一边是大量需要人做的事，一边是没有支付能力对应的工作。人不是没有用处，而是没有被编入任何一个可以付钱给他的结构。\cite{baumol1967macro}

\subhead{真空的形状}

真空不是空无。一个社会不会因为没有人做某件事而出现缺口，因为那些事已经由系统承担了；出现缺口的是人对自己位置的判断。位置被压缩的人不会停在原地等待一个解释，因为没有一个解释可以给出他不知道的事实。他会去找一个仍然能够提供完成感的地方，而这类地方恰好是最不缺的。\key{技能的通胀不降低能力，却取消了能力作为凭据的资格；真空不在需求一侧，而在位置一侧。}

这里说的平庸需要与能力评价分开。平庸指的是一种位置状态：所做的事可以被标准化，因此可被低成本复制；它不描述天赋，也不描述努力的程度。被压缩得最严重的位置，往往由最守纪律、最按时交付、最愿意配合流程的人占据，因为标准化正是在这些人手里被完成的。这也决定了这类压缩很难靠个人努力来避免：在一个可以被标准化的位置上再加努力，只是让这个位置更快地被系统学走。

说到底，真空与平庸是同一件事的两面。表面上，产出变得更容易，交付物的平均质量在上升，一个人可以用更少的时间完成更多的事；结构上，这些产出不再携带关于人的信息。当一个社会里大多数可交付物都可由系统生成时，可交付物就不再能回答“你是谁”这个问题。答案没有消失，只是不能再从产出这一侧取得。寻找新的答案有两条路：一条是退回能够提供完成感的系统，用低成本的替代品填补位置的空缺；另一条是在产出之外重新找到不可外包的位置。前一条路成本低、反馈快、随时可用，因此它先被走到。

哪一条路会被走，取决于位置还有什么别的东西可以挂靠。当功能层的连接断掉，而新的连接还没有建立时，人会先接受任何能提供反馈的东西。这就把问题从功能层推到了真实层。
